%************************************************
\chapter{Background}\label{ch:background}
%************************************************

Graphs in mathematics are used to model pairwise relationships between objects.
Objects are modeled as nodes $\mathcal{V}$, and their relations are called edges
$\mathcal{E}$. Many domains benefit from modeling data as objects (nodes) and
relationships (edges). Consider, for example, social networks, protein--drug
interactions, or railway networks. Studying graph structure and node features allows
us to predict properties of the graph. In fraud detection, for instance, nodes can
represent accounts and edges transactions. By modeling typical graph evolution---such
as which transactions and accounts are added---it becomes possible to detect anomalous
transactions that deviate from expected patterns.

When making predictions about graphs it is often beneficial to first create an
embedding of the graph. The optimal embedding depends on the downstream task, and
there are several ways to produce one: transductively with traditional \acsfont{GRL}
approaches, or inductively with \acsfont{GNN}s.

This chapter provides a fundamental understanding of traditional \acsfont{GRL}
algorithms and neural networks, and describes how \acsfont{GNN}s combine ideas from
both. From traditional \acsfont{GRL} they inherit the idea that node embeddings should
reflect graph structure, while from neural networks they adopt a learnable,
parameterized function. We also discuss challenges that arise when working with
\acsfont{GNN}s on large datasets---challenges that do not exist in traditional neural
networks. Finally, we give a brief introduction to \texttt{PyTorch Geometric}, the
machine learning library used during experiments, and to \texttt{Neo4j}, the graph
database used as graph storage.


%------------------------------------------------
\section{Traditional Graph Representation Learning}\label{sec:traditional_grl}
%------------------------------------------------

In graph representation learning we aim to find a representation of the graph that
encodes its structure. Such embeddings are typically used in downstream tasks such as
node classification, link prediction, or clustering. Finding a good embedding is
somewhat dependent on the downstream task.

The simplest way of constructing node embeddings is to manually select graph
statistics as features. Common examples include node degree, clustering coefficient,
and centrality measures such as eigenvector centrality or PageRank, which aim to
capture local and global importance.

A more advanced class of traditional \acsfont{GRL} methods learns shallow node
embeddings by capturing similarities between nodes. Some of these methods can be
formulated as matrix factorization approaches, where the goal is to approximate a
node--node similarity matrix, often derived from the adjacency matrix or its powers. A
key limitation of manually designed statistics is that they are not learned but
instead depend on what the user believes is important for the task, making them
inflexible and limited in expressive power.

A third type of traditional node embedding is the random-walk-based algorithm. These
methods work by sampling random sequences of nodes and then learn embeddings by
maximizing co-occurrence between nodes that appear within the same random walk. Such
methods have also been shown to perform a form of matrix factorization over
node co-occurrence derived from random walks.

The goal is to learn a decoding function that gives the probability of visiting node
$v$ on a random walk of length $T$ starting from node $u$:
\[
\mathbf{DEC}(\mathbf{z}_u,\mathbf{z}_v)
\triangleq
\frac{e^{\mathbf{z}_u^{\top}\mathbf{z}_v}}
{\sum_{v_k\in\mathcal{V}} e^{\mathbf{z}_u^{\top}\mathbf{z}_k}}
\approx p_{\mathcal{G},T}(v\mid u).
\]


%------------------------------------------------
\subsection{Transductive Nature of Traditional \acsfont{GRL}}\label{sec:transductive}
%------------------------------------------------

A common feature of all \acsfont{GRL} algorithms discussed above---whether using node
statistics, matrix factorization, or random walks---is that no weights are shared
between nodes: no embedding-generating function is learned, only node embeddings for
the current nodes in the graph. If new nodes are introduced, not only is there no way
to compute their embeddings, the embeddings of the existing nodes also become
inaccurate (the latter also occurs if nodes are removed).

Consider a simplified version of the optimization problem for random walk algorithms:
\[
\min_{\{z_x \,\,: \,\,x \in \mathcal{V} \}}
\sum_{(u,v)\in\mathcal{D}}
-\log\!\left(\sigma(\mathbf{z}_u^\top \mathbf{z}_v)\right),
\]
where $\mathcal{D}$ is the set of node pairs encountered on random walks, and
$\mathcal{V}$ is the set of nodes in the graph. The optimization problem above
optimizes the actual node embeddings directly, rather than some embedding-generating
function. This is what makes all traditional representation learning algorithms
\emph{transductive} in nature.

Consider a social network in which people are nodes and their relationships are edges.
If a transductive graph representation is used, we learn feature vectors for specific
people in the graph. If more people are introduced to the network, we must recompute
every person's feature vector.


%------------------------------------------------
\section{Neural Networks}\label{sec:neural_networks}
%------------------------------------------------

In contrast to the transductive nature of traditional \acsfont{GRL}, neural networks
are \emph{inductive}: they learn a generalizable function.

Neural networks are models composed of interconnected processing units called neurons.
How neurons are connected depends on the model architecture. Some neurons receive
external input and pass it further into the network; others aggregate information from
preceding neurons; and output neurons produce the final result. Every connection
between two neurons is associated with a weight, by which the activation of the
upstream neuron is scaled before reaching the next neuron. The network learns patterns
by adjusting these weights.

The neural network can be defined as a function of weights $\theta$ that takes an
arbitrary input $x$ and outputs $\hat{y}$:
\[
f(x \,, \,\theta) = \hat{y}.
\]

The network learns by optimizing $\theta$ to minimize a loss function $\mathcal{L}$
over training data $X$, which quantifies the prediction error---either supervised
(given a desired output $y$) or unsupervised (without a ground-truth outcome):
\begin{equation}
    \min_{\theta} \, \mathcal{L}(X, \theta) = \sum_{x \in X} \, l \, (\,f(x, \theta)\,, \, y \,).
    \tag{supervised}
\end{equation}

The training data consists of input datapoints $\{x \in X_{\text{train}}\}$ mapped to
targets $\{y \in Y_{\text{train}}\}$, forming the training set
$\{(x_1,y_1),(x_2,y_2),\ldots,(x_n,y_n)\}$. The idea is to maximize the likelihood
$p(y_i \mid x_i,\theta)$ of the network producing the correct output for each input by
adjusting the weights, in the hope that the learned weights generalize to unseen data.
The model's generalization ability is evaluated on a separate test set drawn from the
same distribution as the training data.


%------------------------------------------------
\subsection{Weight Optimization}\label{sec:optimization}
%------------------------------------------------

The optimization method used to find the optimal weights is typically some variant of
gradient descent (\acsfont{GD}). \acsfont{GD} involves iteratively updating the model
weights based on the current gradient of the loss function. Differentiating the loss
with respect to the weights $\nabla_{\theta} \mathcal{L}$ gives the direction along
which the loss increases most steeply; hence the negative gradient indicates the
direction of steepest descent.\footnote{%
    \acsfont{GD} converges to the global minimizer for convex differentiable functions
    with appropriate step sizes. The loss function for neural networks is generally
    not convex, but variants of gradient descent are still used to find ``good''
    weights.}
The weights are updated in proportion to the gradient with step size $\eta$, commonly
called the learning rate:
\[
\theta^{(k+1)} = \theta^{(k)} - \eta \, \cdot \nabla_{\theta^{(k)}} \mathcal{L}.
\]

Datapoints are typically assumed to be independent in standard supervised learning,
meaning each prediction $\hat{y}_i$ depends only on $x_i$. Under this assumption the
full gradient is a sum over all datapoints, making vanilla gradient descent
computationally expensive. Stochastic gradient descent (\acsfont{SGD}) approximates
this gradient using a uniformly sampled mini-batch, reducing computation per update.
Although \acsfont{SGD} does not guarantee a decrease in loss at every step, its
expected gradient is an unbiased estimate of the full-batch gradient, yielding
convergence in expectation under the independence assumption.

This approach also alleviates memory constraints. Full-batch gradient descent requires
access to the entire training set for each update, which may not fit in memory. By
computing gradients on mini-batches, \acsfont{SGD} limits memory usage and allows the
batch size to be tuned to the available hardware.


%------------------------------------------------
\section{Graph Neural Networks}\label{sec:gnns}
%------------------------------------------------

Graph neural networks incorporate ideas from traditional graph representation learning
by preserving the intuition that node representations should reflect the graph
structure---\ie, nodes that are close or structurally similar should have similar
embeddings. However, instead of relying on fixed similarity measures or directly
optimizing node embeddings, \acsfont{GNN}s use a neural network architecture to learn
this relationship.

By combining this structural inductive bias with a parameterized, trainable function,
\acsfont{GNN}s learn to generate node embeddings through neighborhood aggregation.
This allows the model to adapt representations to the downstream task via
gradient-based optimization, rather than relying on manually designed features or
predefined objectives. As a result, \acsfont{GNN}s are more expressive and flexible
than traditional \acsfont{GRL} methods. At the same time, they have a fundamentally
different architecture: instead of learning embeddings tied to specific nodes in a
fixed graph, they learn a shared function that generates embeddings for any node based
on its features and local neighborhood.

Returning to the social network example from \autoref{sec:traditional_grl}: when a
new person joins the network, an embedding can be generated directly using the learned
embedding function. Unlike transductive methods, there is no need to retrain the model
or recompute embeddings for existing nodes. The same function can also be applied to
an entirely new social network, as long as the underlying data distribution is similar
to that seen during training. This ability to generalize to unseen nodes and graphs is
what makes \acsfont{GNN}s \emph{inductive}.


%------------------------------------------------
\subsection{Graph Neural Network Architecture}\label{sec:gnn_arch}
%------------------------------------------------

The assumption of independence between datapoints is fundamentally violated in
graph-structured data. In a graph, each node is explicitly defined through its
relationships to other nodes; considering nodes in isolation therefore discards
essential information. The label or representation of a node is often highly dependent
on its neighborhood structure and on the attributes of nearby nodes.

\acsfont{GNN}s address this limitation by explicitly modeling relational dependencies.
Instead of learning a function where the output for a node depends solely on its own
features,
\[
f(x_i, W) = \hat{y}_i,
\]
\acsfont{GNN}s learn representations that incorporate information from a node's local
neighborhood, allowing predictions to depend on both node attributes and the structure
of connections.

\acsfont{GNN}s are built from stacked message-passing layers. At each layer $k$, a
node $v$ updates its representation by aggregating information from its neighbors:
\[
\mathbf{x}_v^{(k)} =
\phi^{(k)}\!\left(
\mathbf{x}_v^{(k-1)},
\bigoplus_{u \in \mathcal{N}(v)}
\psi^{(k)}\!\left(\mathbf{x}_v^{(k-1)}, \mathbf{x}_u^{(k-1)}\right)
\right),
\]
where $\psi^{(k)}$ is a function aggregating a node's neighborhood, and $\phi^{(k)}$
is an update function combining the node's own features with the aggregated
neighborhood information. The aggregation function is typically chosen to be
permutation-invariant with respect to the ordering of neighboring nodes, and the
update function is usually the learnable component.

After $K$ layers, the final node representation encodes information not only about the
node itself but also about its $K$-hop neighborhood. This iterative message passing
enables the \acsfont{GNN} to capture long-range dependencies between $K$-hop
neighbors, which are inaccessible to standard neural networks.

The key inductive bias of \acsfont{GNN}s is therefore relational: predictions are not
made in isolation but are conditioned on the local graph structure. While this
increases computational complexity and introduces challenges for scaling, it allows the
model to exploit information that would otherwise be lost under the \emph{i.i.d.}
assumption.


%------------------------------------------------
\section{Graph Neural Networks and Large Graphs}\label{sec:gnns_large}
%------------------------------------------------

\begin{itemize}
    \item Single layer $\to$ unbiased estimator of mean. Multiple layers $\to$ nested
    sampling $\to$ nonlinear composition $\to$ biased.
    \item GraphSAGE does not correct for unequal sampling probabilities or
    degree-dependent sampling imbalance.
    \item \acsfont{GNN}s often have few layers; that is why they work well in practice
    even though the sampling may be biased.
\end{itemize}

In typical neural networks the datapoints are either completely independent, or in
some cases dependent but with a structure that makes it straightforward to determine
which datapoints depend on which.

\acsfont{GNN}s lack any such relational inductive bias because there are no
architectural assumptions about the dependency structure. Other neural architectures
that exhibit inter-datapoint dependencies usually have some built-in relational
assumption. Time-series models assume a sequential ordering between datapoints, and
convolutional neural networks (\acsfont{CNN}s) assume spatial proximity between
pixels. These assumptions make it straightforward to identify the set of datapoints on
which a given datapoint depends.

In a time-series model one can easily derive that datapoint $t_n$ depends on all
$\{t_{n-1}, t_{n-2},\ldots,t_{n-k}\}$. In a \acsfont{CNN}, pixel $p_{x,y}$ depends
only on
\[
\{\,p_{i,j} \mid |i-x| \le r,\ |j-y| \le r \,\},
\]
where $r$ is the receptive-field radius, and pixels are never dependent on pixels in
other images.

We say that $x_1$ depends on $x_2$ if:
\begin{equation}
    \label{eq:dep}
    f(x_1) \neq f(x_1 \mid x_2)
    \quad \text{or} \quad
    f(x_1) \neq f(x_1, x_2).
\end{equation}

There are different degrees of strictness in such dependencies. Time-series models
often maintain an internal state summarizing past observations: datapoints only need
to be presented in the correct order. To make an accurate prediction at timestep $t$
one merely needs to have previously inserted $x_0, x_1, \ldots, x_{t-1}$ and then
insert $x_t$ at the current step (the relaxed form in the left-hand side of
\autoref{eq:dep}). \acsfont{CNN}s and \acsfont{GNN}s are stricter: all datapoints that
$x$ depends on must be provided as immediate inputs to the model.


%------------------------------------------------
\subsection{Training on Large Datasets Under Memory Constraints}\label{sec:memory}
%------------------------------------------------

Training on a dataset of $X$\,GB on a machine with $Y$\,GB of memory where $Y < X$
is straightforward for typical neural networks with \acsfont{SGD}, because it suffices
to load datapoints in mini-batches
$\{d^{(1)}, d^{(2)}, \ldots, d^{(n)}\}$ such that
$\sum_{i} d^{(i)} < Y$. For weakly structured dependencies one additionally needs to
load datapoints in the correct sequence; for immediate structural dependencies (as in
\acsfont{CNN}s) one must ensure correct batch partitioning.

The problem is compounded for \acsfont{GNN}s. The output for node $x$ depends on its
$K$-hop neighborhood. If the average node degree is $n$, then the number of nodes
each node directly depends on is approximately $n^K$.

Returning to the social network fraud-detection example: predicting whether a person
will make a fraudulent transaction requires, for a 2-layer \acsfont{GNN} with an
average of 50 transactions per account, approximately $50 \times 50 = 2{,}500$ times
more datapoints than a traditional neural network that uses only that person's
features.


%------------------------------------------------
\subsection{Subgraph Sampling}\label{sec:subgraph_sampling}
%------------------------------------------------

The node-to-node dependency means that after sampling a mini-batch of seed nodes, one
must build a subgraph containing the seed nodes and their neighbors.

The $K$-hop neighborhood of seed nodes can vary considerably in size. To avoid
exploding neighborhood sizes, a sampling algorithm is typically used to select a
subset of the true $K$-hop neighborhood. Sampling strategies vary in implementation,
but they all impose an upper bound on the number of nodes sampled at each hop. These
bounds are usually specified as $[k_0, k_1, \ldots, k_n]$, where $k_i$ is the maximum
number of neighbors to sample at hop $i$ from the seed nodes.

This uniform neighbor sampling renders the mini-batch gradient biased, and therefore
does not guarantee convergence for \acsfont{SGD}. In practice, however, convergence
is generally sufficient:
\begin{equation*}
\mathbb{E}_S\!\left[
\nabla_W \ell\!\left(
\sigma\!\left(
\frac{1}{|\mathcal{N}^S(v)|}
\sum_{u\in\mathcal{N}^S(v)} h_u^{k-1}
\right)
\right)
\right]
\;\neq\;
\nabla_W \ell\!\left(
\sigma\!\left(
\frac{1}{|\mathcal{N}(v)|}
\sum_{u\in\mathcal{N}(v)} h_u^{k-1}
\right)
\right),
\end{equation*}
where the left-hand side is the expected mini-batch gradient over sampled neighbor
graphs, and the right-hand side is the full gradient over all neighbors.

Previous work has addressed this problem by partitioning the graph across several
worker machines that each train the \acsfont{GNN} using the gradient computed from
their partition. Partitions are constructed so that clusters of nodes with many common
neighbors reside on the same machine, thereby minimizing inter-worker communication.

The paper ``Scalable and Efficient Full-Graph \acsfont{GNN} Training for Large
Graphs'' implements a system along these lines with the graph partitioned among $k$
workers. Each worker iteratively generates the output for each message-passing layer,
synchronized with the other workers; after each layer it both sends its computed
hidden states to workers that require them and receives hidden states for nodes stored
elsewhere.

Similarly, the paper ``ByteGNN: Efficient Graph Neural Network Training at Large
Scale'' partitions the graph among a number of workers. The high-level architecture
consists of one \texttt{GraphStore} responsible for partitioning, paired with
sampling and training workers. The sampling workers draw a subgraph around the
seed nodes in the mini-batch and fetch required neighboring nodes from other workers.

Finally, Neo4j has presented a paper ``Graph Neural Networks on Graph Databases''
demonstrating that it is possible to train a \acsfont{GNN} by storing the graph in a
Neo4j graph database and sampling subgraphs around the seed nodes for each mini-batch.
This is, to the best of our knowledge, the only scientific paper exploring a
graph-database solution to the problem of scaling graph neural networks.


%------------------------------------------------
\section{PyTorch Geometric}\label{sec:pyg}
%------------------------------------------------

PyTorch Geometric is a graph neural network library that enables easy creation of
\acsfont{GNN}s in Python. It is built on top of PyTorch and provides helpful modules
such as predefined message-passing layers (graph attention, graph convolution, and
more), making it straightforward to construct customizable \acsfont{GNN}s by combining
different layers within implementations of the abstract PyTorch \texttt{nn.Module}
class.

PyTorch Geometric also provides a framework for storing the graph remotely in a
database, through the following abstract classes:

\begin{itemize}
    \item \texttt{FeatureStore:} Defines how node feature vectors are stored and
    fetched. A fast feature-vector look-up is essential.

    \item \texttt{GraphStore:} Defines how the graph topology is stored. In contrast
    to the \texttt{FeatureStore}, the emphasis here is on fast traversal rather than
    fast look-up, since efficient neighbor sampling requires rapid graph traversal.

    \item \texttt{BaseSampler:} Allows the user to define a sampling method for
    selecting neighbors used in the forward pass. The \texttt{BaseSampler} is tightly
    coupled with the \texttt{GraphStore}. Multiple sampling processes can be spawned;
    since instances of the sampler are recreated in each separate process, it is
    important not to store large amounts of data in the sampler.

    \item \texttt{NodeLoader:} Takes implementations of \texttt{FeatureStore},
    \texttt{GraphStore}, and \texttt{BaseSampler} together with batching
    hyperparameters (\eg, batch size) and provides a data loader that splits the data
    into mini-batches, samples a subgraph for each batch, and fetches the
    corresponding node feature vectors from the \texttt{FeatureStore}.
\end{itemize}


%------------------------------------------------
\section{Neo4j}\label{sec:neo4j}
%------------------------------------------------
